\documentclass[10pt]{report}
\usepackage{epsf}
\usepackage{amsmath}
\usepackage{amssymb}
\usepackage{palatino}
\usepackage[pdftex]{graphics}
\usepackage{fancyhdr}
\usepackage{epsfig}
\usepackage{multirow}
\usepackage{multicol}
\usepackage{cancel}
\usepackage{hyperref}
\usepackage{longtable}
\usepackage{ulem}
\usepackage{draftwatermark}
\usepackage{etoolbox}
\usepackage{enumitem}
\setlist[description]{leftmargin=10pt,labelindent=10pt}

\parindent 0in
\parskip 1ex
\oddsidemargin  0in
\evensidemargin 0in
\textheight 8.5in
\textwidth 6.5in
\topmargin -0.25in

\pagestyle{fancy}
\newcommand{\leftheader}{MedTech Prototyping Skills (BME290L)}
\newcommand{\rightheader}{Palmeri (Fall 2023)}
\lhead{\textbf{\leftheader}}
\rhead{\textbf{\rightheader}}
\lfoot{Compiled: \today}
\cfoot{}
\rfoot{\thepage}

\newcommand{\Holiday}[2]{%
    \options{#1}{\noclassday}
    \caltext{#1}{#2 (No Class)}
}

\begin{document}

\section*{Class Syllabus}

\subsection*{Personnel}
\begin{itemize}
    \item[] \textbf{Instructor:} Dr. Mark Palmeri
    \begin{itemize}
        \item Email: \href{mailto:mark.palmeri@duke.edu}{mark.palmeri@duke.edu}\footnote{Response times will be way
        faster via Teams.}
        \item Office Hours: \url{http://meet.palmeri.io}
        \begin{itemize}
            \item 258 Hudson Hall Annex
            \item \url{https://duke.zoom.us/my/mark.palmeri}
        \end{itemize}
    \end{itemize}
    \item[] \textbf{Lab Master:} Matt Brown \href{mailto:matt.brown@duke.edu}{(matt.brown@duke.edu)}
    \item[] \textbf{Teaching Assistants:} 
    \begin{itemize}
        \item Katy Price \href{mailto:catherine.price@duke.edu}{(catherine.price@duke.edu)}
    \end{itemize}
\end{itemize}

\subsection*{Course Times \& Locations}
\begin{itemize}
    \item[] \textbf{Lecture:} Mon \& Wed from 10:05--11:20 in Wilkinson 136

    \item[] \textbf{Lab:} Fri from 08:30--11:30 in Wilkinson 126 /
    Teer212\footnote{This lab is available 24/7 to you, expect when other
    classes have the lab reserved.} (Door Code: 5-34-1)
\end{itemize}

\subsection*{Course Objectives}
This course focuses on developing medical device prototyping skills that will be
used in future project and design courses, with a focus on preparing our
students for positions in the medtech industry.  Students will work individually
to complete design tasks that will be tested to quantitative specifications.
Students will gain hands-on experience with device fabrication, debugging,
testing and failure analysis.

Upon completion of this course, students should be able to:

\begin{itemize}
    \item Perform functional decomposition and develop user action flow chart
    \item Utilize ECAD (KiCad) for:
    \begin{itemize}
        \item Electronic schematic capture
        \item Printed circuit board (PCB) layout
        %\item SPICE circuit simulation
    \end{itemize}
    \item Modular breadboarding of circuits to translation to testable PCBs
    \item Design a battery-powered device that:
    \begin{itemize}
        \item Accepts analog and digital user input
        \item Outputs analog and digital outputs
        \item Performs logic using discrete electronics (e.g., logic gates, 555 timers)
        \item Passes testable specifications to 95\% CI
    \end{itemize}
    \item Implement electronic logic on microcontroller (Arduino framework)
    \begin{itemize}
        \item Modular / testable code development in C
        \item Software version control (\texttt{git})
        \item Interupt Service Routines
        \item Pulse Width Modulation
    \end{itemize}
    \item Utilize CAD (Onshape) for:
    \begin{itemize}
        \item Device enclosure design with UI/UX considerations
        \item Input / Output considerations
        \item Size / weight constraints
        \item Preparation of mechanical drawings
    \end{itemize}
    \item 3D printing of designs
    \begin{itemize}
        \item Assembly of discrete enclosure pieces
        \item Use of mechanical fasteners
    \end{itemize}
    \item Write technical analysis documents / reports (Jupyter notebooks)
\end{itemize}


\subsection*{Prerequisites}
\begin{itemize}
    \item Introduction to Design (EGR101)
    \item Introductory Circuit Analysis (ECE110 or equivalent, corequisite)
\end{itemize}

\subsection*{Learning Management System}
We will be using \href{https://canvas.duke.edu/}{Canvas}--not Sakai--this
semester, and it will be homebase for all things course related, including:
\begin{itemize}
    \item Ed Discussion - the fastest way to get questions answered (including
    anonymously) from Dr. Palmeri and the teaching assistants
    \item Resources, including:
    \begin{itemize}
        \item Electronics Textbooks
        \item Lab Equipment Manuals \& Training Tutorials
        \item Online Software \& Tools
        \item Required Hardware to Purchase
    \end{itemize}
    \item All course content modules, assignments, quizzes, and lab exercises.
    \item Due dates for everything.
    \item Grades
\end{itemize}

Some assignments / quizzes / labs will be submitted via
\href{https://www.gradescope.com/}{Gradescope}, which is integrated into Canvas.
When applicable, \textbf{you must associate the pages of your submission with
the grading rubric criteria during the submission process.}  Failure to do so
will result in lost assignment credit.

\subsection*{Resources}
\subsection*{Hardware}
One of the following development kits is \textbf{required} hardware for this semester:
\begin{itemize}
    \item \href{https://store-usa.arduino.cc/products/arduino-nano-33-ble?selectedStore=us}{Arduino Nano 33 BLE} 
    \item \href{https://store-usa.arduino.cc/products/arduino-nano-33-ble-sense?selectedStore=us}{Arduino Nano 33 Sense}
    \item \href{https://store-usa.arduino.cc/products/arduino-nano-every-with-headers?selectedStore=us}{Arduino
    Nano Every}\footnote{This board is being used in BME354L this semester.}
\end{itemize}

\subsubsection*{Textbooks}
All of these books contain valuable content and will be referenced throughout
the semester.  If you are looking to pursue a career in medical device design,
then it may be worth having one of these as a reference.

\begin{itemize}
    \item \href{https://find.library.duke.edu/catalog/DUKE009381868}{Practical Electronics for Inventors (Scherz \& Monk) [Fourth Edition]}\footnote{This is available online through the Duke Library.}
\end{itemize}

\subsubsection*{Online Resources}
We will be using the following [E]CAD packages:
\begin{itemize}
    \item \href{https://duke.onshape.com}{Onshape}\footnote{Note, this is a
    specific Duke instance of Onshape, not the public one.  You will be invited
    to join by Dr. Palmeri.}
    \footnote{If you know Fusion 360 or SolidWorks and would rather use those CAD
    packages, you are welcome to, but the class will only support Onshape.}
    \item \href{https://kicad.org/}{KiCad}
    \footnote{Please download and install the stable version (currently \texttt{v6.0.10}). Do not install the
    \texttt{v7.x} release candidate!}
    \footnote{If you know Altium Designer, you are welcome to use it instead, but
    the class will only support KiCad.}
\end{itemize}

These are cloud tools or cross-platform for installation on your personal computers.

\begin{itemize}
    \item \href{https://teams.microsoft.com}{Teams} - ``Slack-like'' tool for live
    chat, team discussion, asking questions and resource sharing
    \item \href{https://diagrams.net}{diagrams.net} - software / hardware flowcharts, functional decomposition
    \item \href{https://git-scm.com/}{Git} - version control software
    \item \href{https://gitlab.oit.duke.edu}{GitLab} - version control / project mangement
    \item One of the following for firmware development:
    \begin{itemize}
        \item \href{https://code.visualstudio.com/}{Visual Studio Code} with \href{https://platformio.org/}{PlatformIO} extension
        \item \href{https://www.arduino.cc/en/software}{Arduino IDE}
    \end{itemize}
\end{itemize}

\subsection*{Attendance \& Participation}
Class participation, including lecture attendance and scheduled lab time,
contributes to your class grade.   Not being able to participate in class
activities due to illness should be reported using the
\href{http://www.pratt.duke.edu/undergrad/policies/3531}{Short Term Illness Form
(STIF)} \textbf{before} the missed class activity.

\textbf{If you have a non-illness-related reason (e.g., a job interview) for not
being able participate in a class activity or meet an assignment submission
deadline, please reach out to Dr.\ Palmeri as soon as possible to discuss the
situation.}

Students are responsible for obtaining missed lecture content from other
students in the class.  Lectures will be recorded via Panopto and available via
Canvas.

\subsection*{Class Schedule} 
The class is organized in a sequence of modules.  Specific details surrounding
dates for assignments, quizzes, labs, etc. will be detailed on
Canvas/Gradescope.

This course uses a version of
\href{https://en.wikipedia.org/wiki/Mastery_learning}{mastery learning}, where
``mastery'' of a given module is necessary to progress onto the subsequent
module.  In this course, the assignments/quizzes/labs of one module must be
completed before moving onto the next module.

\subsubsection*{Modules}
\begin{enumerate}
    \item Fundamentals \& Engineering Tools
    \item Computer Aided Design (CAD) [OnShape]
    \item ECAD: Schematic Capture (KiCad)
    \item ECAD: PCB Layout [KiCad]
    \item Version Control [Git]
    \item Firmware Development (Arduino Framework / PlatformIO)
    \item Breadboard Testing / Electronics Bench Equipment
    \item PCB Fabrication
    \item Enclosure Design
    \item Integration \& Testing
\end{enumerate}

\subsubsection*{Special Dates of Note}
\begin{description}
    \item[October 16, 2023] No Class (Fall Break)
    \item[November 22, 2023] No Class (Thanksgiving)
    \item[December 08, 2023 at 17:00] Final Project Due (absolutely no late submissions accepted)
\end{description}

\subsection*{Grading} 
The following grading scheme will be used for this course:

\begin{center}
\begin{tabular}{ll}
Participation \& Quizzes             & 25\% \\
Assignments \& Lab Technical Reports & 50\% \\
Final Project Report                 & 25\% \\
\end{tabular}
\end{center}

\subsubsection*{Final Course Grades}
This course is not ``curved'' (i.e., a distribution of grades will not be
enforced), and a traditional grading scheme will be used (e.g., 90-93 = A-,
94-97 = A, 97-100 = A+).

Failing the course can happen with a cummulative score $<$ 65 or not completing
all of the assignments.  \textbf{All assignments must be satisfactorily
completed by each student to pass the class, even if no credit will be awarded
based on the late policy.}

All grades will be posted to Gradescope/Canvas throughout the semester to
track your performance.


\subsubsection*{Late Policy}
Permission to submit an assignment late should be sought from Dr. Palmeri as far
in advance as reasonably possible, but no less than 48 hours in advance, except
in cases of acute illness.

Unexcused late assignments will lose 50\% of their potential point value for
each 24 hour period beyond the due date (i.e., 100\% \textrightarrow 50\% →
25\%\ldots), but remember that all assignments must be completed at a
satisfatory level to pass the course, even if no grade credit is assigned.

\textbf{No late final projects will be accepted!  Not submitting a final project
on time will prevent passing the class.}

\subsubsection*{Regrades}
Any regrading requests need to be made via Gradescope within one week of grades
for a given assignment being returned using Gradescope. You must provide a
description of why you feel a regrade is appropriate. Requesting a regrade could
lead to additional loss of credit when an assignment is re-evaluated.

Some assignments will have an opportunity to be resubmitted based on grading
feedback at the discretion of Dr. Palmeri.

\subsection*{Collaboration, Citing Sources, Duke Community Standard \& Academic Honor}
\subsubsection*{Collaboration} Engineering is inherently a collaborative field,
and in this class, you are encouraged to work collaboratively on your projects.
That being said, all of the work that you submit must be generated by you and
reflect your understanding of the material. All resources developed by another
person or company, and used in your project, must be properly acknowledged.

\subsubsection{Duke Community Standard}
All students are expected to adhere to all principles of the
\href{https://trinity.duke.edu/undergraduate/academic-policies/community-standard-student-conduct}{Duke
Community Standard}. Violations of the Duke Community Standard will be referred
immediately to the Office of Student Conduct. Please do not hesitate to talk
with Dr. Palmeri about any situations involving academic honor, especially if it
is ambiguous what should be done.

\subsubsection*{Can I use AI?}
The use of artificial intelligence is a rapidly developing resource / tool in
engineering.  In software development, there are many levels of AI-assitance
available.  Such form of assistance include the
\href{https://visualstudio.microsoft.com/services/intellicode/}{IntelliCode}
tools, \href{https://github.com/features/copilot}{GitHub CoPilot}, and Large
Language Models (LLMs).  Similar to using ``traditional'' resources like
\href{https://stackoverflow.com/}{Stack Overflow} for programming guidance,
these tools can be leveraged to help with syntax.  You are, however, strongly
cautioned to not rely on these tools for logical implementation.  \textbf{Any
use of AI-assisted tools should be acknowledged in your submission through
comments in your code.}

\end{document}
